% --- Archon physics formalization source begin ---
% archon:physics
% archon:covers IPhO2026Problems/problem_IPhO_2026_3_B_1.lean
% archon:source-report reports/ipho_2026_k3/problem_IPhO_2026_3_B_1.source.json
% archon:problem-id IPhO_2026_3
% archon:part-id B.1
% archon:previous-part IPhO_2026_3_A_3
% archon:previous-part-policy natural_language_prerequisite_only

\chapter{Physics problem IPhO\_2026\_3\_B\_1}
\label{ch:IPhO2026Problems_problem_IPhO_2026_3_B_1}

\paragraph{Problem source.}
Continue with the paramagnetic torus.  Its equation of state is T*M*V = n*K*H,
its heat capacity at constant M is C\_M = n*lambda/T\textasciicircum{}2, and dU = C\_M*dT.
The volume is fixed and the magnetic work on the material is
dW = mu\_0*V*H*dM.  Work and heat entering the torus are positive.

Current subquestion:
At fixed temperature T, H changes from H\_i to H\_f. Find the heat Q transferred into the torus.

\paragraph{Current subquestion.}
At fixed temperature T, H changes from H\_i to H\_f. Find the heat Q transferred into the torus.

\paragraph{Recorded answer/context.}
Q = -(mu\_0*n*K/(2*T))*(H\_f\textasciicircum{}2 - H\_i\textasciicircum{}2).

\paragraph{Figure/image path.}
/root/proposal\_for\_physic/science-mango/ipho\_2026\_source/image/T3\_page-3.png

\paragraph{Reusable previous-part conclusions.}
\begin{itemize}
\item Source A.3. Question: Subtract the vacuum-core contribution and write the work dW done on the paramagnetic material. Reusable conclusions: dW = mu\_0*V*H*dM. Policy: natural\_language\_prerequisite\_only; do\_not\_import\_Lean\_output
\end{itemize}

\paragraph{Formalization target.}
create a compiling Lean file with sorry bodies at `IPhO2026Problems/problem\_IPhO\_2026\_3\_B\_1.lean`.
The Lean declarations must preserve the physical quantities, dimensions or dimensional roles, figure labels, governing-law hypotheses, and final relation expressed by this problem.
Use Mathlib/Physlib names found through LeanExplore where available. If a domain API is missing, introduce faithful local abstractions rather than scalar placeholder aliases.

\begin{theorem}[Physics formalization target]
\label{thm:physics:IPhO_2026_3_B_1:target}
\uses{thm:IPhO2026Problems_problem_IPhO_2026_3_B_1:OfficialAnswerValue}
The assigned autoformalize agent should translate this physics problem into Lean declarations in the covered file, with theorem and lemma proof bodies written as `by sorry`.
\end{theorem}
\begin{proof}
This is an autoformalization task, not a proof task. Produce faithful statements that can later be proved without weakening the source contract.
\end{proof}
% NOTE: PhysLean-coverage exemption (planner-recorded, iter-002): PhysLean's thermodynamics modules do not cover this paramagnetic-torus magnetic-work/isothermal-first-law model (see this file's physics-grounding log). Self-containment is kept with the `import Mathlib` baseline; no irrelevant Physlib import is added. This documents the resolved import policy (iter-001 review ruling) for the blueprint-doctor `missing-physlib-import` check.
% --- Archon physics formalization source end ---

% --- Archon named-quantities coverage (blueprint-writer 3-b-1-entries) ---
% NOTE: import policy (mirrors the reconciliation note above, iter-002 exemption):
% the covered file builds on the `import Mathlib` baseline alone; PhysLean has no
% paramagnetic-torus magnetic-work or isothermal first-law module, so no Physlib
% import is added.  This NOTE records the resolved import policy for the chapter.
\subsection*{Quantities and data}

\begin{definition}[Torus parameters]
\label{def:IPhO2026Problems_problem_IPhO_2026_3_B_1:TorusParams}
\lean{IPhO2026.Problem3.B1.TorusParams}
The constant parameters of the paramagnetic torus experiment: the
permeability of free space $\mu_0$, the fixed torus volume $V$, the amount
of paramagnetic material $n$ (moles), the material constant $K$ appearing
in the equation of state $T M V = n K H$, and the heat-capacity coefficient
$\lambda$ of $C_M = n\lambda/T^2$.  All fields are physical scalars;
positivity and regularity constraints enter separately as hypotheses where
they are needed.
\end{definition}
\begin{proof}
Definition; no claim.
\end{proof}

\begin{definition}[Torus state]
\label{def:IPhO2026Problems_problem_IPhO_2026_3_B_1:TorusState}
\lean{IPhO2026.Problem3.B1.TorusState}
The variables that vary along a quasistatic process of the torus: the
temperature $T$, the magnetization $M$ (magnetic moment per unit volume),
and the applied magnetic field strength $H$.  The volume is fixed, so it
lives in the parameters record rather than here.
\end{definition}
\begin{proof}
Definition; no claim.
\end{proof}

\begin{definition}[Heat capacity at constant magnetization]
\label{def:IPhO2026Problems_problem_IPhO_2026_3_B_1:HeatCapacityConstM}
\lean{IPhO2026.Problem3.B1.heatCapacityConstM}
\uses{def:IPhO2026Problems_problem_IPhO_2026_3_B_1:TorusParams}
The heat capacity at constant magnetization as given in the problem
statement, $C_M(T) = n\lambda/T^2$, as a real-valued readout of the torus
parameters at temperature $T$.
\end{definition}
\begin{proof}
Definition; no claim.
\end{proof}

\subsection*{Governing-law structures}

\begin{definition}[Equation of state of the paramagnet]
\label{def:IPhO2026Problems_problem_IPhO_2026_3_B_1:SatisfiesEOS}
\lean{IPhO2026.Problem3.B1.SatisfiesEOS}
\uses{def:IPhO2026Problems_problem_IPhO_2026_3_B_1:TorusParams, def:IPhO2026Problems_problem_IPhO_2026_3_B_1:TorusState}
The magnetic equation of state of the torus material as a predicate on a
parameter record and a state: $T\,M\,V = n\,K\,H$.  This is a governing law
recorded from the problem statement, not a target.
\end{definition}
\begin{proof}
Definition; no claim.
\end{proof}

\begin{definition}[Heat-capacity law for the internal energy]
\label{def:IPhO2026Problems_problem_IPhO_2026_3_B_1:HasHeatCapacityLaw}
\lean{IPhO2026.Problem3.B1.HasHeatCapacityLaw}
\uses{def:IPhO2026Problems_problem_IPhO_2026_3_B_1:TorusParams, def:IPhO2026Problems_problem_IPhO_2026_3_B_1:HeatCapacityConstM}
The characterization of the material's internal energy $U$: at constant
magnetization $dU = C_M\,dT$; that is, $U$ is a function of temperature
alone, differentiable at every nonzero $T$ with derivative
$C_M(T) = n\lambda/T^2$.
\end{definition}
\begin{proof}
Definition; no claim.
\end{proof}

\begin{definition}[Magnetic work-on density]
\label{def:IPhO2026Problems_problem_IPhO_2026_3_B_1:IsMagneticWorkDensity}
\lean{IPhO2026.Problem3.B1.IsMagneticWorkDensity}
\uses{def:IPhO2026Problems_problem_IPhO_2026_3_B_1:TorusParams}
The pointwise record of the magnetic work differential of part A.3,
$dW_{\mathrm{on}} = \mu_0\,V\,H\,dM$, along a quasistatic process
parametrized by the magnetization: a density function satisfies
$\mathit{workDensity}(M) = \mu_0\,V\,H(M)$ for every $M$, so the work done
on the material from $M_0$ to $M_1$ is
$\int_{M_0}^{M_1} \mathit{workDensity}(M)\,dM$.  Natural-language
prerequisite of part A.3, recorded as a governing law.
\end{definition}
\begin{proof}
Definition; no claim.
\end{proof}

\begin{definition}[First law with magnetic work, isothermal form]
\label{def:IPhO2026Problems_problem_IPhO_2026_3_B_1:ObeysFirstLawMagnetic}
\lean{IPhO2026.Problem3.B1.ObeysFirstLawMagnetic}
\lean{IPhO2026.Problem3.B1.IsothermalFieldChange}
\lean{IPhO2026.Problem3.B1.IsothermalFieldChange.field_increases}
\lean{IPhO2026.Problem3.B1.IsothermalFieldChange.h_branch}
\uses{def:IPhO2026Problems_problem_IPhO_2026_3_B_1:TorusParams, def:IPhO2026Problems_problem_IPhO_2026_3_B_1:SatisfiesEOS, def:IPhO2026Problems_problem_IPhO_2026_3_B_1:HasHeatCapacityLaw, def:IPhO2026Problems_problem_IPhO_2026_3_B_1:IsMagneticWorkDensity}
The first law of thermodynamics for the torus along a quasistatic
isothermal process at temperature $T_{\mathrm{iso}}$, under the source sign
convention (work and heat entering the torus are positive): given
$\mathit{HasHeatCapacityLaw}$, the internal-energy bracket vanishes over
every isothermal leg, so for all $M_0, M_{\mathrm{target}}$,
\[
Q_{\mathrm{in}}(M_{\mathrm{target}}) - Q_{\mathrm{in}}(M_0)
= \bigl(U(T_{\mathrm{iso}}) - U(T_{\mathrm{iso}})\bigr)
- \int_{M_0}^{M_{\mathrm{target}}} \mathit{heatDensity}(M)\,dM .
\]
The accompanying \texttt{IsothermalFieldChange} record packages the process
data of one isothermal field change $H_i \to H_f$ at fixed $T \neq 0$: the
magnetization branch $H \mapsto M(H)$ obeying the equation of state on the
whole tracked range (the convex hull of $H_i$, $H_f$, and the demagnetized
reference field $0$), the cumulative heat readout $Q_{\mathrm{in}}(M)$
normalized by $Q_{\mathrm{in}}(0) = 0$, and a Boolean orientation flag for
the field ramp certified against the endpoint magnetizations.
\end{definition}
\begin{proof}
Definition; no claim.  The first law yields the leg balance by
$\Delta U = 0$ along an isotherm; the field-change record only packages
data and certifications.
\end{proof}

\subsection*{Derivation bridges}

\begin{lemma}[Magnetization fixed by the equation of state]
\label{lem:IPhO2026Problems_problem_IPhO_2026_3_B_1:MagnetizationEqEosSolution}
\lean{IPhO2026.Problem3.B1.magnetization_eq_eos_solution}
\uses{def:IPhO2026Problems_problem_IPhO_2026_3_B_1:SatisfiesEOS, def:IPhO2026Problems_problem_IPhO_2026_3_B_1:ObeysFirstLawMagnetic}
Along the tracked isotherm of an \texttt{IsothermalFieldChange} with
$V \neq 0$, the pointwise solution of the equation of state pins the
magnetization to a linear function of the applied field,
$M(H) = n K H / (T V)$, at every $H$ in the recorded range.
\end{lemma}
\begin{proof}
Solve $T M V = n K H$ for $M$ using $T \neq 0$ and $V \neq 0$: divide
$T M V = n K H$ through by $T V$.
\end{proof}

\begin{lemma}[Leg stays in the tracked range]
\label{lem:IPhO2026Problems_problem_IPhO_2026_3_B_1:LegMemTrackedRange}
\lean{IPhO2026.Problem3.B1.leg_mem_tracked_range}
\uses{def:IPhO2026Problems_problem_IPhO_2026_3_B_1:ObeysFirstLawMagnetic, lem:IPhO2026Problems_problem_IPhO_2026_3_B_1:MagnetizationEqEosSolution}
The magnetization-ramp leg from the demagnetized reference $0$ to any
tracked endpoint magnetization $M(H)$ stays inside the tracked field range,
so the equation of state --- hence
\cref{lem:IPhO2026Problems_problem_IPhO_2026_3_B_1:MagnetizationEqEosSolution}
--- applies to every point of the work integrals
$\int_{0}^{M(H)} \mathit{workDensity}(M)\,dM$.
\end{lemma}
\begin{proof}
Case analysis on the position of $M \in [0, M(H)]$ against the endpoint
membership bounds; both range endpoints are $\min$/$\max$ combinations, so
bound chaining closes the pure set-inclusion obligation --- physical
content is only the convexity of the tracked range.
\end{proof}

\begin{lemma}[Evaluation of one isothermal work leg]
\label{lem:IPhO2026Problems_problem_IPhO_2026_3_B_1:LegWorkIntegralEval}
\lean{IPhO2026.Problem3.B1.leg_work_integral_eval}
\uses{def:IPhO2026Problems_problem_IPhO_2026_3_B_1:IsMagneticWorkDensity, def:IPhO2026Problems_problem_IPhO_2026_3_B_1:ObeysFirstLawMagnetic, lem:IPhO2026Problems_problem_IPhO_2026_3_B_1:MagnetizationEqEosSolution, lem:IPhO2026Problems_problem_IPhO_2026_3_B_1:LegMemTrackedRange}
The magnetic work delivered on the torus along the isothermal ramp from the
demagnetized reference to a tracked field $H$ evaluates to
\[
\int_{0}^{M(H)} \mu_0 V M(H(M))\,dM
= \frac{\mu_0 n K}{2 T}\,\bigl(M(H)\bigr)^2 ,
\]
the micro-evaluation used once per endpoint in the target proof.
\end{lemma}
\begin{proof}
Rewrite the integrand pointwise via
\cref{lem:IPhO2026Problems_problem_IPhO_2026_3_B_1:MagnetizationEqEosSolution}
(applicable on the whole leg by
\cref{lem:IPhO2026Problems_problem_IPhO_2026_3_B_1:LegMemTrackedRange}),
so the integrand becomes $\mu_0 V \cdot M$ times a constant scalar; then
evaluate $\int c\,M\,dM = c\,M(H)^2/2$.
\end{proof}

\subsection*{Isothermal heat and value theorems}

\begin{definition}[Heat transferred into the torus]
\label{def:IPhO2026Problems_problem_IPhO_2026_3_B_1:HeatTransferredIntoTorus}
\lean{IPhO2026.Problem3.B1.heatTransferredIntoTorus}
\uses{def:IPhO2026Problems_problem_IPhO_2026_3_B_1:ObeysFirstLawMagnetic}
The physical heat transferred into the torus between the initial and final
states of the isothermal field change: the difference of the tracked
cumulative heat readout,
$Q_{\mathrm{in}}(M(H_f)) - Q_{\mathrm{in}}(M(H_i))$.  This is the quantity
the target theorem characterizes; its closed form must still be derived.
\end{definition}
\begin{proof}
Definition; no claim.
\end{proof}

\begin{definition}[Closed-form heat value carrier]
\label{def:IPhO2026Problems_problem_IPhO_2026_3_B_1:HeatIntoTorusValue}
\lean{IPhO2026.Problem3.B1.heat_into_torus_value}
\uses{def:IPhO2026Problems_problem_IPhO_2026_3_B_1:TorusParams, def:IPhO2026Problems_problem_IPhO_2026_3_B_1:ObeysFirstLawMagnetic}
The named closed-form value of the heat transferred into the torus along an
isothermal field change, $-(\mu_0 n K/(2T))(H_f^2 - H_i^2)$.  This only
\emph{names} the quantity the target theorem speaks about; no hypothesis
states that the physical heat equals this value --- deriving the equality
from the equation of state, the work law, and the first law is the proof
obligation.
\end{definition}
\begin{proof}
Definition; no claim.
\end{proof}

\begin{theorem}[Isothermal heat into the torus]
\label{thm:IPhO2026Problems_problem_IPhO_2026_3_B_1:IsothermalHeatIntoTorus}
\lean{IPhO2026.Problem3.B1.isothermal_heat_into_torus}
\uses{def:IPhO2026Problems_problem_IPhO_2026_3_B_1:TorusParams, def:IPhO2026Problems_problem_IPhO_2026_3_B_1:HasHeatCapacityLaw, def:IPhO2026Problems_problem_IPhO_2026_3_B_1:IsMagneticWorkDensity, def:IPhO2026Problems_problem_IPhO_2026_3_B_1:ObeysFirstLawMagnetic, def:IPhO2026Problems_problem_IPhO_2026_3_B_1:HeatTransferredIntoTorus, def:IPhO2026Problems_problem_IPhO_2026_3_B_1:HeatIntoTorusValue, lem:IPhO2026Problems_problem_IPhO_2026_3_B_1:MagnetizationEqEosSolution, lem:IPhO2026Problems_problem_IPhO_2026_3_B_1:LegWorkIntegralEval}
Along an isothermal field change of the paramagnetic torus obeying the
equation of state, with the magnetic work $dW = \mu_0 V H\,dM$ of part A.3
and the first law of thermodynamics, the heat transferred into the torus
(the difference of the tracked heat readout between the final and initial
states) equals the closed-form value carrier:
$Q = -(\mu_0 n K/(2T))(H_f^2 - H_i^2)$.  No hypothesis states this closed
form in advance: the first-law hypothesis only supplies leg balances
against the demagnetized reference, and deriving the value of each work
integral is the proof obligation.
\end{theorem}
\begin{proof}
Apply the first-law leg balance to the two legs $0 \to M(H_i)$ and
$0 \to M(H_f)$; the internal-energy brackets vanish by the heat-capacity
law along the isotherm, so
$Q = -\bigl(W_{\mathrm{on}}(H_f) - W_{\mathrm{on}}(H_i)\bigr)$ with each
$W_{\mathrm{on}}$ evaluated by
\cref{lem:IPhO2026Problems_problem_IPhO_2026_3_B_1:LegWorkIntegralEval};
substituting the equation-of-state solution
$M(H) = n K H /(T V)$ at both endpoints and cancelling $V$ reduces the
difference to the closed form by algebra.
\end{proof}

\begin{theorem}[Recorded official answer]
\label{thm:IPhO2026Problems_problem_IPhO_2026_3_B_1:OfficialAnswerValue}
\lean{IPhO2026.Problem3.B1.official_answer_value}
\uses{def:IPhO2026Problems_problem_IPhO_2026_3_B_1:TorusParams, def:IPhO2026Problems_problem_IPhO_2026_3_B_1:ObeysFirstLawMagnetic, def:IPhO2026Problems_problem_IPhO_2026_3_B_1:HeatIntoTorusValue}
The closed value of the heat transferred into the torus carried by the
target theorem equals the closed form recorded in the official answer key,
\[
Q = -\frac{\mu_0\,n\,K}{2 T}\,\bigl(H_f^2 - H_i^2\bigr).
\]
This entry is conclusion-side only: it pins the value of the answer
carrier; it is not a hypothesis of any theorem.
\end{theorem}
\begin{proof}
Definitional unfolding of the value carrier; reflexivity.
\end{proof}
